\section{Introduction}\label{sec:introduction}
We consider trace estimation when a matrix can be accessed only through
matrix--vector products. Under a fixed query budget, low-rank variance
reduction creates an allocation problem: a query used to learn a basis is
no longer available to sample the residual. The relevant question is not
simply whether a matrix has a rapidly decaying spectrum. It is whether the
realized basis reduces the residual estimator risk enough to compensate for
the queries spent constructing it.

\subsection{Background and problem formulation}
\hutchpp\ combines a randomized low-rank trace contribution with stochastic
estimation of the complementary trace. Meyer et al.~\cite{meyer2021}
establish improved matrix--vector query complexity for positive semidefinite
matrices. Our starting point is this decomposition, rather than a competing
claim of improved worst-case query complexity.

In a basic implementation, $q$ queries form a range sketch, $r$ queries
evaluate the accepted rank-$r$ basis, and $\ell$ queries estimate the residual.
Thus $q+r+\ell=m$. The familiar denominator $m-2q$ is a full-rank special
case, not the general ledger. Moreover, the equalities $r=q$ and
$q=r_\star$, where $r_\star$ is a dominant spectral rank, do not imply that the
randomized basis accurately captures that dominant eigenspace.

This distinction leads to two related questions. First, what exact
conditional-risk reduction justifies another low-rank direction? Second,
can this reduction be certified from a small number of additional queries?
The second question includes an economic constraint: confidence information
uses queries, and an abstaining procedure cannot recover those queries by
returning to a baseline basis.

\subsection{Results and scope}
We organize the analysis around four contributions.
\begin{enumerate}
\item \textbf{Exact allocation and marginal-risk identities.}
We state the conditions under which random, sequentially chosen allocations
remain unbiased. We distinguish Gaussian and Rademacher conditional risk and
show that an accepted direction is beneficial precisely when its fractional
residual-energy reduction exceeds the fractional loss of residual capacity
(Proposition~\ref{thm:T4}).
\item \textbf{Capture mechanisms beyond numerical rank.}
For step spectra, a nonnegative residual-energy decomposition separates tail
energy from signal leakage. A graph representation explains the effect of a
poorly conditioned square signal sketch. A separate exact construction shows
that a leading eigendirection can be absent from an otherwise full-rank pilot
(Proposition~\ref{thm:T6}).
\item \textbf{Certification targets and cost.}
We derive the moments and Hoeffding decomposition of sample variance,
formulate paired risk differences, and give safe-acceptance implications
under simultaneous confidence events. We distinguish those implications
from complete-policy safety after construction and certification costs
(Propositions~\ref{thm:T9} and~\ref{thm:T10}).
\item \textbf{Explicit limitations supported by frozen experiments.}
Two audited confidence constructions are analytically vacuous on the
small-budget grid. Empirical studies nevertheless find useful risk
information, substantial common-probe covariance reduction, and rare-path
protection. We report the accompanying failures and limitations rather than
interpreting these observations as a universal guarantee.
\end{enumerate}

Many underlying identities are elementary or classical. We do not claim
them individually as new mathematical discoveries. Their role is to make the
allocation decision, the probe-specific risk target, and the cost of
certification explicit in a common framework. The contribution of the
experimental analysis is a controlled separation of failure mechanisms,
including hypotheses that the recorded results did not support.

\subsection{Relation to prior work}
\paragraph{Randomized range finding.}
Halko, Martinsson, and Tropp~\cite{halko2011} develop a general framework for
randomized low-rank approximation. We use a self-contained step-spectrum
calculation to study capture, rather than importing a Gaussian sketch
theorem for a rotated coordinate-Rademacher block.

\paragraph{Quadratic-form concentration.}
Cortinovis and Kressner~\cite{cortinovis2022} provide the quadratic-chaos tail
bound used in our truncation calculation. Hypercontractivity supplies
degree-dependent moment control~\cite{odonnell2007}. Neither ingredient,
without the additional steps analyzed below, is a useful small-budget
certificate for the realized risk comparison.

\paragraph{Organization and boundary.}
We follow a theory-first exposition: preliminaries, allocation analysis,
capture mechanisms, certification, and experimental validation, with
supporting calculations and provenance in appendices. The detailed
experimental records are unpublished project artifacts~\cite{project2026}.
All empirical claims are conditional on their stated finite populations.
We do not claim a universal one-column oversampling rule, a new minimax lower
bound, or a completed confidence-certified online allocator.
